\chapter{Signal Models}

\section{The contract}

Every signal implements \texttt{generate(close) -\/> DataFrame} of scores, same
shape as \texttt{close}. A score is a ``desire to hold'' in roughly $[-1, 1]$;
the portfolio layer (Chapter~6) turns scores into leverage- and risk-constrained
weights. A signal marks itself \texttt{market\_neutral = True} when it already
produces cross-sectionally demeaned scores.

\section{Rule-based signals}

\begin{table}[h]
\footnotesize
\centering
\begin{tabularx}{\textwidth}{L{3.4cm} X}
\toprule
\thh{Signal} & \thh{Score} \\
\midrule
\texttt{Momentum} & $\tanh\!\big(\mathrm{mom}_n / \overline{|\mathrm{mom}_n|}_{252}\big)$, optionally divided by realised vol so noisier names are sized down \\
\texttt{MeanReversion} & $-z_n / (2\cdot\text{entry\_z})$, clipped to $[-1,1]$: fade a stretched short-horizon z-score \\
\texttt{CrossSectionalMomentum} & cross-sectional z-score of $\mathrm{mom}_n$, renormalised each day so $\sum_s |w_s| \approx 1$: long the relative winners, short the relative losers, market-neutral \\
\bottomrule
\end{tabularx}
\end{table}

\texttt{CrossSectionalMomentum} is the pipeline's default and the one used for
the headline demo result -- it is the closest rule-based analogue of the
persistent per-symbol $\alpha_t$ injected into the synthetic data (Chapter~2),
so it is the fairest baseline for the ML comparison in Chapter~8.

\section{A walk-forward ML signal}

\texttt{MLSignal} wraps a \texttt{scikit-learn} pipeline
(\texttt{StandardScaler} + \texttt{HistGradientBoostingRegressor} by default) and
trains it \emph{only} on the walk-forward splits described in Chapter~5: for
each \texttt{(train\_start, train\_end, test\_start, test\_end)} window it fits
on the training rows of the leak-free dataset from Chapter~3, predicts on the
test rows, and stitches the out-of-sample predictions back into a wide panel.
The cross-sectional z-score of that prediction becomes the score, exactly as in
\texttt{CrossSectionalMomentum} -- so the two signals are compared on equal
footing in Chapter~8: same sizing, same costs, same risk limits, different
alpha source.

\section{An optional sequence model}

\texttt{SequenceSignal} is a single-layer LSTM over a rolling window of the same
feature tokens, trained walk-forward with the same split logic:

\begin{figure}[h]
\centering
\begin{tikzpicture}[node distance=9mm]
  \node[stage] (in) {window of \\ $F$ features $\times$ $w$ days};
  \node[stage, right=of in] (lstm) {LSTM \\ hidden size $h$};
  \node[stage, right=of lstm] (head) {linear head};
  \node[stage, right=of head] (out) {forward-return \\ prediction};
  \draw[flowarr] (in)--(lstm); \draw[flowarr] (lstm)--(head); \draw[flowarr] (head)--(out);
\end{tikzpicture}
\caption{One (symbol, date) training example: the last $w$ days of the feature
panel for that symbol, target = forward return over \texttt{horizon} days. The
final hidden state feeds a linear head; only the last time step's output is
used (\texttt{out[:, -1, :]}).}
\end{figure}

It requires \texttt{torch} and is import-guarded: without the \texttt{ml} extra
installed, selecting \texttt{signal.type: sequence} raises a clear
\texttt{ImportError} naming the alternatives, rather than crashing deep inside a
training loop.

\begin{caveat}
This is included to \emph{learn the architecture and benchmark it honestly},
not because it is expected to win. Chapter~10 explains why an LSTM over
technical features is dominated by transaction costs and regime shift on real
data far more often than it earns its keep. Never trust a sequence model's
in-sample equity curve.
\end{caveat}
